\chapter{Configuring Kernel Creation}\label{r01:ch:kernel-creation}

The same symbolic description can produce very different kernels. Whether you get a
scalar loop nest, an OpenMP-parallel AVX kernel or a CUDA function is decided by the
\term{configuration} you pass to the code generator. This chapter describes the entry
point \code{create_kernel}, the configuration object, and the options that matter most.
\Cref{r01:app:config-reference} lists every option.

\section{The entry point}

\begin{pycode}
def create_kernel(assignments: SymbolicKernel,
                  config: CreateKernelConfig | None = None,
                  **kwargs) -> Kernel
\end{pycode}
\code{assignments} may be a single \code{Assignment}, a list of assignments, an
\code{AssignmentCollection}, or a flowgraph (\cref{r01:sec:flow}). The configuration can be
given in three ways:
\begin{pycode}
ps.create_kernel(asms)                                   # all defaults
ps.create_kernel(asms, cfg)                              # a CreateKernelConfig object
ps.create_kernel(asms, cfg, target=ps.Target.X86_AVX)   # cfg with fields replaced
\end{pycode}
Keyword arguments replace \emph{top-level} fields of the configuration
(\code{dataclasses.replace}). Nested options such as \code{cpu.openmp.enable} cannot be
passed this way, so for those you need a configuration object. The function lives in
\file{src/pystencils/codegen/driver.py}. It creates a
\cls{DefaultKernelCreationDriver} and calls it. \Cref{r01:ch:architecture} describes what
the driver does.

\section{The configuration object}\label{r01:sec:config}

\cls{CreateKernelConfig} (\file{codegen/config.py}) is a dataclass of
\term{options}, grouped into \term{categories}:
\begin{pycode}
cfg = ps.CreateKernelConfig(target=ps.Target.X86_AVX, default_dtype="float32")
cfg.ghost_layers = 1                        # top-level option
cfg.cpu.openmp.enable = True                # category cpu, sub-category openmp
cfg.cpu.vectorize.enable = True
cfg.gpu.indexing_scheme = "blockwise4d"     # ignored for CPU targets
\end{pycode}
Each option is a descriptor whose value is \code{None} while it is \emph{unset}. Reading
an unset option returns its default. \code{cfg.get_option(name)} and
\code{cfg.is_option_set(name)} make this distinction explicit. Validators normalize
input, so \code{default_dtype="float32"} is converted to a \code{PsIeeeFloatType} and the
string \code{"blockwise4d"} to \code{GpuIndexingScheme.Blockwise4D}. Two methods help
when you work with several variants:
\begin{itemize}
  \item \code{cfg.copy()} makes an independent (semi-deep) copy, so changing a copy never
        affects the original:
        \pysnippet{ex07_vectorize_x86.copy}
  \item \code{cfg.override(other)} copies all options that are \emph{set} in
        \code{other} into \code{cfg}. Frameworks use this to layer a user's
        configuration on top of their own defaults.
\end{itemize}

\begin{table}[htbp]
  \centering\small
  \caption{The most important options (full list in \cref{r01:app:config-reference}).}
  \label{r01:tab:config-overview}
  \begin{tabularx}{\linewidth}{@{}l l >{\raggedright\arraybackslash}X@{}}
    \toprule
    \textbf{Option} & \textbf{Default} & \textbf{Meaning}\\
    \midrule
    \code{target} & \code{GenericCPU} & hardware target (\cref{r01:sec:targets})\\
    \code{default_dtype} & \code{float64} & type of untyped symbols and dynamic fields\\
    \code{index_dtype} & \code{int64} & type of counters, sizes and strides\\
    \code{ghost_layers} & \code{AUTO} & iteration space: skip cells at the border\\
    \code{iteration_slice} & -- & iteration space: an explicit slice\\
    \code{index_field} & -- & iteration space: sparse list of cells\\
    \code{function_name} & \code{"kernel"} & name of the generated function\\
    \code{jit} & per target & JIT compiler used by \code{Kernel.compile()}\\
    \code{cpu.openmp.enable} & \code{False} & parallelize the outermost loop\\
    \code{cpu.vectorize.enable} & \code{False} & explicit SIMD vectorization\\
    \code{cpu.vectorize.assume_inner_stride_one} & \code{False} & contiguous fastest coordinate\\
    \code{gpu.indexing_scheme} & \code{Linear3D} & mapping of cells to GPU threads\\
    \code{gpu.manual_launch_grid} & \code{False} & user sets block and grid size\\
    \bottomrule
  \end{tabularx}
\end{table}

\section{Targets}\label{r01:sec:targets}

\cls{Target} (\file{codegen/target.py}) is a \code{Flag} enum. Each public target is
the combination of private feature bits (\code{_CPU}, \code{_VECTOR}, \code{_X86},
\code{_AVX}, \dots), and the predicates \code{is_cpu()}, \code{is_vector_cpu()},
\code{is_gpu()} and \code{is_automatic()} test those bits.

\begin{center}\small
\begin{tabularx}{\linewidth}{@{}l l >{\raggedright\arraybackslash}X@{}}
  \toprule
  \textbf{Target} & \textbf{Kind} & \textbf{Notes}\\
  \midrule
  \code{GenericCPU} (\code{CPU}) & CPU & portable C, OpenMP possible, no SIMD\\
  \code{CurrentCPU} & CPU, automatic & resolved to \code{Target.auto_cpu()}: the best vector target of the host\\
  \code{X86_SSE}, \code{X86_AVX}, \code{X86_AVX512}, \code{X86_AVX512_FP16} & vector CPU & 128/256/512-bit, x86-64-v2/v3/v4\\
  \code{ARM_NEON}, \code{ARM_NEON_FP16} & vector CPU & 128-bit\\
  \code{ARM_SVE} & vector CPU & scalable; lane count must be set explicitly\\
  \code{RISCV_RVV} & vector CPU & \emph{fork only}: VLEN = 256 (\cref{r01:ch:rvv-backend})\\
  \code{CUDA}, \code{HIP} & GPU & run through CuPy\\
  \code{CurrentGPU} (\code{GPU}) & GPU, automatic & CUDA or HIP, depending on the installed CuPy\\
  \code{SYCL} & SYCL & experimental; run through dpctl\\
  \bottomrule
\end{tabularx}
\end{center}

The capabilities of the host are detected with py-cpuinfo on x86 and ARM Linux,
\code{sysctlbyname} on macOS, and \file{/proc/cpuinfo} on RISC-V. On the laptop used to
generate this manual:
\outsnippet{ex07_vectorize_x86.targets}
The first entry is \code{GenericCPU}. Its canonical flag name is the private bit
\code{_CPU}, which is how it prints. \code{default_vector_lanes(dtype)} is the number of
lanes the vectorizer uses when \code{cpu.vectorize.lanes} is unset: the register width
divided by the element width.

\section{Data types}\label{r01:sec:data-types}

Two options control the types the code generator chooses on its own:
\begin{description}
  \item[\code{default_dtype}] (\code{float64}): applied to free untyped symbols, to
        fields with the dynamic type \code{ps::numeric_t}, and to
        \code{DynamicType.NUMERIC_TYPE} symbols.
  \item[\code{index_dtype}] (\code{int64}): the type of loop counters, sizes, strides
        and \code{DynamicType.INDEX_TYPE} symbols. When index values are themselves
        vectorized, as for the Philox counters of a vectorized random number generator,
        the index type must have the width of the vector elements. A \code{float32}
        Philox kernel for \code{X86_AVX} fails with \code{int64} (``x86/AVX does not
        support int64<8>'') and works with \code{index_dtype="int32"}.
\end{description}
A generic field declared as \code{"f: [2D]"} therefore produces \code{double} or
\code{float} code depending on \code{default_dtype}. Fields with an explicit type keep
that type, and mixing types requires explicit casts (\cref{r01:sec:types-frontend}).

\section{The iteration space}\label{r01:sec:iteration-space}


\begin{figure}[htbp]
  \centering
  % Cells visited (blue) by the kernels of Section 4.5 on an 8 x 8 array.
% x = coordinate 0 (ctr_0) to the right, y = coordinate 1 (ctr_1) upwards.
\begin{tikzpicture}[x=0.34cm,y=0.34cm,font=\scriptsize]
  \newcommand{\ispanel}[4]{% x-shift, filter (TeX condition on \x,\y), title, code
    \begin{scope}[shift={(#1,0)}]
      \foreach \x in {0,...,7} \foreach \y in {0,...,7} {
        \pgfmathparse{#2 ? 1 : 0}
        \ifnum\pgfmathresult=1 \fill[psblue!50] (\x,\y) rectangle ++(1,1); \fi
      }
      \draw[gridline] (0,0) grid (8,8);
      \draw[domainline,line width=0.5pt] (0,0) rectangle (8,8);
      \node[align=center,font=\footnotesize] at (4,9.2) {#3};
      \node[align=center,font=\ttfamily\tiny] at (4,-1.5) {#4};
    \end{scope}}
  \ispanel{0}{(\x>=1) && (\x<=6) && (\y>=1) && (\y<=6)}
    {(a) ghost layers}{ghost\_layers=1\\(AUTO for Jacobi)}
  \ispanel{10.5}{(\x>=1) && (\x<=4)}
    {(b) per side}{ghost\_layers=\\{}[(1, 3), 0]}
  \ispanel{21}{(mod(\x,2)==0) && (\y>=3) && (\y<=4)}
    {(c) slice}{iteration\_slice=\\make\_slice[::2, 3:-3]}
  \ispanel{31.5}{(\y>=\x)}
    {(d) triangle}{iteration\_slice=\\make\_slice[:, ctr\_0:]}
  \draw[axisarrow] (0,-3.4) -- (3,-3.4) node[right] {$x$ (\texttt{ctr\_0})};
  \draw[axisarrow] (-0.6,0) -- (-0.6,3) node[above] {$y$};
\end{tikzpicture}

  \caption[Iteration-space options]{Cells visited (blue) by kernels with different iteration-space options on an $8\times8$ array. The options are explained below; (a)--(d) correspond to the kernels of this section.}
  \label{r01:fig:ispaces}
\end{figure}

The iteration space is the set of cells the kernel visits (\cref{r01:fig:ispaces}). It is defined by \emph{at
most one} of three options. Setting two of them raises a \code{ValueError}.

\subsection{Ghost layers}
For dense kernels the default is \code{ghost_layers=AUTO}. The code generator takes the
largest neighbour offset over all accesses to domain fields and skips that many cells on
every side. An explicit number, or a \code{(low, high)} pair per axis, overrides this:
\pysnippet{ex05_iteration_space.ghost-layers}
\csnippet[nolol=false,caption={\texttt{ghost\_layers=[(1, 3), 0]}: one cell at the low end and three at the high end of axis 0.}]{ex05_iteration_space.gl-sides}
The chosen value is stored in the kernel's metadata, so that a framework can check that
its arrays have enough ghost layers:
\outsnippet{ex05_iteration_space.gl-metadata}

\subsection{Iteration slices}
\code{iteration_slice} restricts the kernel to a slice, written with
\code{ps.make_slice}. Bounds may be negative (relative to the end), and steps are
allowed. With a fixed-size field all bounds and strides become constants:
\pysnippet{ex05_iteration_space.slice}
\csnippet{ex05_iteration_space.slice}
Slice bounds may even depend on the counters of \emph{slower} dimensions, which gives
non-rectangular iteration spaces such as a triangle:
\pysnippet{ex05_iteration_space.triangle}
\csnippet{ex05_iteration_space.triangle}

\subsection{Sparse kernels with index fields}\label{r01:sec:sparse}
For boundary conditions, or any other update on an irregular set of cells, a kernel can
iterate over a list of cell coordinates instead of a dense box. The list is a 1D field of
type \code{INDEXED} with a struct element type. Its members \code{x}, \code{y}
(\code{z}) are the coordinates, and any further members are available as
\code{idx("member")} (\cref{r01:fig:sparse}):
\pysnippet{ex05_iteration_space.sparse}
\outsnippet{ex05_iteration_space.sparse-out}
\csnippet[nolol=false,caption={A sparse kernel. The struct type was lowered to byte offsets.}]{ex05_iteration_space.sparse}
This is how \code{pystencils.boundaries} implements boundary handling.

\begin{figure}[htbp]
  \centering
  % The sparse kernel of Section 4.5.3: one loop iteration per entry of the index list.
\begin{tikzpicture}[font=\footnotesize]
  % ---- index list --------------------------------------------------------------
  \matrix (idx) [matrix of nodes, nodes={draw=psgray!70,minimum width=0.95cm,
      minimum height=0.55cm,font=\ttfamily\scriptsize,anchor=center},
      column sep=-\pgflinewidth, row sep=-\pgflinewidth,
      row 1/.style={nodes={fill=psgray!15,font=\scriptsize\bfseries}}]
  {
    x & y & value\\
    0 & 2 & 3.0\\
    1 & 3 & 42.0\\
    2 & 1 & 5.0\\
  };
  \node[above=2pt of idx,align=center] {index field \texttt{idx}\\[-1pt]\scriptsize(\texttt{INDEXED}, struct elements)};
  \foreach \r/\l in {2/0,3/1,4/2} \node[left=2pt of idx-\r-1,font=\ttfamily\scriptsize,text=psgray] {sparse\_idx=\l};

  % ---- domain field ------------------------------------------------------------
  \begin{scope}[shift={(5.0,-1.6)},x=0.85cm,y=0.85cm]
    \fill[psblue!45] (0,2) rectangle ++(1,1);
    \fill[psblue!45] (1,3) rectangle ++(1,1);
    \fill[psblue!45] (2,1) rectangle ++(1,1);
    \draw[gridline] (0,0) grid (3,4);
    \node[celltext] at (0.5,2.5) {3.0};
    \node[celltext] at (1.5,3.5) {42.0};
    \node[celltext] at (2.5,1.5) {5.0};
    \node[above,align=center] at (1.5,4) {domain field \texttt{dom} (3 $\times$ 4)};
    \draw[axisarrow] (0,-0.3) -- (1.2,-0.3) node[right,celltext] {$x$};
    \draw[axisarrow] (-0.3,0) -- (-0.3,0.9) node[left,celltext] {$y$};
    \coordinate (c1) at (0.0,2.5); \coordinate (c2) at (1.0,3.5); \coordinate (c3) at (2.0,1.5);
  \end{scope}
  \draw[flowarrow,psblue] (idx-2-3.east) to[out=0,in=180] (c1);
  \draw[flowarrow,psblue] (idx-3-3.east) to[out=0,in=200] (c2);
  \draw[flowarrow,psblue] (idx-4-3.east) to[out=0,in=180] (c3);
  \node[align=left,anchor=north west,font=\ttfamily\scriptsize] at (-1.6,-2.3)
    {for (sparse\_idx = 0; sparse\_idx < 3; ++sparse\_idx) \{\\
     \ \ ctr\_0 = idx[sparse\_idx].x;\ \ ctr\_1 = idx[sparse\_idx].y;\\
     \ \ dom(ctr\_0, ctr\_1) = idx[sparse\_idx].value;\ \}};
\end{tikzpicture}

  \caption[A sparse index-list kernel]{The sparse kernel above: the loop runs over the entries of the index list, and each entry supplies the coordinates of one cell plus additional data.}
  \label{r01:fig:sparse}
\end{figure}


\section{CPU optimizations}\label{r01:sec:cpu-options}

\paragraph{OpenMP.}
\code{cpu.openmp.enable = True} turns the outermost loop into an OpenMP work-sharing
loop inside a parallel region (see \cref{r01:lst:reduction} for a reduction). The options
\code{schedule} (default \code{"static"}), \code{num_threads} and \code{collapse} map
directly to the corresponding OpenMP clauses.

\paragraph{Vectorization.}
\code{cpu.vectorize.enable = True} makes \ps{} emit explicit SIMD intrinsics for the
target's instruction set. The innermost loop is split into a vectorized main part and a
scalar remainder loop. This works only for vector CPU targets; on \code{GenericCPU} the
vectorizer raises a \code{CodegenError}. The options are:
\begin{itemize}
  \item \code{lanes}: the vector width in elements. By default it is the full register
        (4 doubles for AVX).
  \item \code{assume_inner_stride_one}: replaces the stride of the fastest coordinate by
        the constant 1, so that contiguous vector loads and stores can be used. Without
        it, the stride is a runtime value, and strided accesses need gather and scatter
        instructions. On x86, gathers need AVX2 and scatters AVX-512, so a kernel that
        writes a strided field cannot be vectorized for SSE or AVX
        (\cref{r01:lst:matrix}).
\end{itemize}
\pysnippet{ex07_vectorize_x86.config}
\csnippet[nolol=false,caption={The Jacobi kernel for \texttt{X86\_AVX} with OpenMP and vectorization.},label=r01:lst:avx]{ex07_vectorize_x86.code}

\begin{figure}[htbp]
  \centering
  % Loop structure of the AVX Jacobi kernel (Listing 4.x) for a 30 x 37 array.
\begin{tikzpicture}[font=\footnotesize]
  % ---- outer loop: OpenMP -------------------------------------------------------
  \begin{scope}[x=0.44cm,y=0.44cm]
    \colorlet{th0}{psteal!45}\colorlet{th1}{psblue!35}\colorlet{th2}{psorange!45}\colorlet{th3}{psred!25}
    \fill[ghostcell] (0,0) rectangle ++(1,1);
    \fill[ghostcell] (29,0) rectangle ++(1,1);
    \foreach \t/\c in {0/th0,1/th1,2/th2,3/th3} \fill[\c] (1+7*\t,0) rectangle ++(7,1);
    \draw[gridline] (0,0) grid (30,1);
    \foreach \t/\c in {0/th0,1/th1,2/th2,3/th3}
      \node[celltext,fill=\c,inner sep=1pt] at (4.5+7*\t,0.5) {thread \t};
    \node[anchor=south west,font=\footnotesize] at (0,1.15)
      {\textbf{outer loop} \texttt{ctr\_0 = 1 \dots{} 28}: \texttt{\#pragma omp for schedule(static)}, e.g.\ 4 threads};
    \node[celltext,text=psgray] at (0.5,-0.45) {0};
    \node[celltext,text=psgray] at (29.5,-0.45) {29};
  \end{scope}

  % ---- inner loop of one row: SIMD blocks + remainder ---------------------------
  \begin{scope}[shift={(0,-1.9)},x=0.357cm,y=0.44cm]
    \fill[ghostcell] (0,0) rectangle ++(1,1);
    \fill[ghostcell] (36,0) rectangle ++(1,1);
    \foreach \b in {0,...,7} {
      \pgfmathsetmacro{\shade}{mod(\b,2)==0 ? 50 : 28}
      \fill[psblue!\shade] (1+4*\b,0) rectangle ++(4,1);
      \draw[black!60,line width=0.6pt] (1+4*\b,0) rectangle ++(4,1);
    }
    \fill[remcell] (33,0) rectangle ++(3,1);
    \draw[gridline] (0,0) grid (37,1);
    \foreach \b in {0,...,7} \draw[black!60,line width=0.6pt] (1+4*\b,0) rectangle ++(4,1);
    \node[anchor=south west,font=\footnotesize] at (0,1.15)
      {\textbf{inner loop} for each \texttt{ctr\_0}: \texttt{ctr\_1 = 1 \dots{} 35} (contiguous)};
    \draw[decorate,decoration={brace,mirror,amplitude=4pt}] (1,-0.15) -- (33,-0.15)
      node[midway,below=5pt,align=center,celltext]
        {8 vector iterations, 4 lanes each: \texttt{\_mm256\_loadu\_pd}, \texttt{\_mm256\_add\_pd}, \dots,
         \texttt{\_mm256\_storeu\_pd}};
    \draw[decorate,decoration={brace,mirror,amplitude=4pt}] (33,-0.15) -- (36,-0.15)
      node[midway,below=5pt,align=center,celltext] {scalar\\remainder};
    \draw[psred,thick] (33,-0.1) -- (33,1.6) node[above,celltext,text=psred]
      {\texttt{ctr\_1\_\_rem\_start} = 33};
    \node[celltext,text=psgray] at (0.5,-0.45) {0};
    \node[celltext,text=psgray] at (36.5,-0.45) {36};
  \end{scope}
\end{tikzpicture}

  \caption[OpenMP and SIMD loop structure]{Loop structure of the kernel in \cref{r01:lst:avx} for a $30\times37$ array: OpenMP distributes the outer loop over threads; within each row, 32 cells are processed in 8 vector iterations of 4 lanes, and the last 3 cells by the scalar remainder loop.}
  \label{r01:fig:vectorization}
\end{figure}

\Cref{r01:fig:vectorization} shows the loop structure of this kernel.
In \cref{r01:lst:avx}, \code{ctr_1__rem_start} marks where the remainder loop begins. The
inner strides are gone from the signature, and the constant $1/4$ exists both as a
broadcast vector and as a scalar. The kernel was compiled and checked against NumPy on
an array whose inner size (35) is not a multiple of 4:
\pysnippet{ex07_vectorize_x86.run}

\begin{pitfall}[\texttt{assume\_inner\_stride\_one} is not checked at run time]
  Once the inner stride has been assumed to be 1, it is no longer a kernel parameter. A
  non-contiguous array (for example a NumPy view \code{a[:, ::2]}) is accepted without
  any error and produces \emph{wrong results}. \Cref{r01:app:pitfalls} shows this case.
  The source contains a \texttt{TODO} to pass the assumption on to the runtime as a
  precondition.
\end{pitfall}

\paragraph{Loop blocking.}
\code{cpu.loop_blocking = (bx, by, ...)} tiles the loop nest for cache reuse, with one
entry per spatial dimension. Block sizes may be integers, symbols (runtime parameters),
or \code{None} to leave a dimension unblocked. The block loops are placed outside the
regular loops; with OpenMP enabled, the outermost block loop is the parallel one:
\pysnippet{ex07_vectorize_x86.blocking}
\csnippet[nolol=false,caption={Loop blocking with $16\times8$ tiles and OpenMP.},basicstyle=\ttfamily\scriptsize]{ex07_vectorize_x86.blocking}

\paragraph{Options without effect.}
The options \code{cpu.vectorize.assume_aligned} and
\code{cpu.vectorize.use_nontemporal_stores} are defined, but the current backend never
reads them. x86 kernels always use unaligned loads and stores. Setting
\code{cpu.use_cacheline_zeroing = True} raises a \code{NotImplementedError}.

\section{GPU options}\label{r01:sec:gpu-options}

For \code{Target.CUDA} and \code{Target.HIP}, \ps{} generates a \code{__global__}
function in which each thread handles one cell (or several, with grid striding):
\pysnippet{ex09_gpu.cuda}
\csnippet[nolol=false,caption={CUDA kernel with the default \texttt{Linear3D} indexing.},label=r01:lst:cuda]{ex09_gpu.linear3d}

\begin{figure}[htbp]
  \centering
  % Linear3D GPU indexing of the Jacobi kernel (Listing 4.x): one thread per cell.
% Horizontal: ctr_1 (fastest coordinate) <-> x;  vertical: ctr_0 <-> y.
\begin{tikzpicture}[x=0.5cm,y=0.5cm,font=\footnotesize]
  % array with one ghost layer: 12 (ctr_1) x 8 (ctr_0) cells, interior 10 x 6
  \fill[ghostcell] (0,0) rectangle (12,8);
  \fill[white] (1,1) rectangle (11,7);
  % threads outside the interior: masked by the guard
  \fill[pattern=north west lines,pattern color=psred!70] (11,1) rectangle (13,9);
  \fill[pattern=north west lines,pattern color=psred!70] (1,7) rectangle (11,9);
  % one highlighted thread
  \fill[psblue!55] (7,2) rectangle ++(1,1);
  \draw[gridline] (0,0) grid (12,8);
  \draw[domainline] (1,1) rectangle (11,7);
  % thread blocks of 4 x 4
  \foreach \bx in {0,1,2} \foreach \by in {0,1}
    \draw[psblue!80!black,line width=1.1pt,dashed] (1+4*\bx,1+4*\by) rectangle ++(4,4);
  \node[celltext,text=psblue!80!black,fill=white,inner sep=1pt] at (7,8.45) {block (1, 1)};
  \node[celltext,text=psblue!80!black,fill=white,inner sep=1pt] at (7,0.45) {block (1, 0)};
  % annotations
  \draw[psgray] (7.5,2.5) -- (13.6,2.5) node[right,align=left,celltext]
    {\texttt{blockIdx = (1, 0)}, \texttt{threadIdx = (2, 1)}\\
     \texttt{ctr\_1 = 1 + 1$\cdot$4 + 2 = 7}\\
     \texttt{ctr\_0 = 1 + 0$\cdot$4 + 1 = 2}};
  \draw[psgray] (12.5,6) -- (13.6,6) node[right,align=left,celltext]
    {threads outside the domain do nothing:\\
     \texttt{if (ctr\_0 < \_size\_dst\_0 - 1 \&\& ctr\_1 < \_size\_dst\_1 - 1)}};
  \draw[axisarrow] (0,-0.6) -- (3,-0.6) node[right,celltext]
    {\texttt{ctr\_1} $\leftrightarrow$ \texttt{blockIdx.x*blockDim.x + threadIdx.x}};
  \draw[axisarrow] (-0.6,0) -- (-0.6,3) node[above,align=center,celltext] {\texttt{ctr\_0}\\$\leftrightarrow$ \texttt{.y}};
\end{tikzpicture}

  \caption[\texttt{Linear3D} GPU indexing]{\texttt{Linear3D} indexing of \cref{r01:lst:cuda} with $4\times4$ thread blocks (dashed): each thread computes its cell from its block and thread indices. The grid is rounded up to whole blocks, so the guard \texttt{if} disables the threads outside the domain (hatched).}
  \label{r01:fig:gpu-mapping}
\end{figure}

The \code{gpu.indexing_scheme} option selects the mapping from cells to threads:
\begin{description}
  \item[\code{Linear3D}] (default, up to 3D): every spatial dimension is mapped to
        \code{blockIdx * blockDim + threadIdx} of one CUDA dimension. The fastest
        coordinate goes to $x$. A guard \code{if} skips threads outside the domain
        (\cref{r01:fig:gpu-mapping}).
  \item[\code{GridstridedLinear3D}]: like \code{Linear3D}, but each thread loops over the
        domain with the grid size as stride, so a small grid can cover a large domain.
  \item[\code{Blockwise4D}] (up to 4D): the slow dimensions go to \code{blockIdx}, the
        fastest to \code{threadIdx.x}, i.e.\ one block per line of cells.
\end{description}
Every \code{GpuKernel} carries a factory for its \term{launch configuration}, which
computes block and grid sizes from the array shapes at call time:
\pysnippet{ex09_gpu.schemes}
\outsnippet{ex09_gpu.schemes}
A \code{DynamicBlockSizeLaunchConfiguration} starts from a default block size and can
adapt it to the iteration space (\code{fit_block_size}, \code{trim_block_size}). With
\code{gpu.manual_launch_grid = True}, the user sets \code{block_size} and
\code{grid_size} explicitly. This is required when the code generator cannot compute
the size of the iteration space, as for the triangle above.
\pysnippet{ex09_gpu.launch}
\outsnippet{ex09_gpu.launch}
The option \code{gpu.assume_warp_aligned_block_size}, combined with an explicitly set
\code{gpu.warp_size}, enables warp-level reductions (\cref{r01:sec:gpu-platforms}).

\section{SYCL}
\code{Target.SYCL} generates a function to be called inside a SYCL parallel
command. With \code{sycl.automatic_block_size=True} (the default) the kernel takes a
\code{sycl::item} and is launched over a \code{sycl::range}. Otherwise it takes a
\code{sycl::nd_item} for \code{nd_range} launches. The SYCL JIT (\code{SYCLJit},
compiler \texttt{icpx} by default) runs such kernels from Python on dpctl USM arrays.
SYCL support is marked experimental.

\section{Constraints and code-generation errors}\label{r01:sec:constraints}

Before generating code, \ps{} checks that the kernel can be executed in any loop order
and in parallel. During generation it checks types and target capabilities. The script
\file{examples/ex12_errors.py} triggers the typical errors:
\outsnippet[style=out,basicstyle=\ttfamily\scriptsize]{ex12_errors.errors}
Note the first case: the independence check only looks at reads that come
\emph{after} a write. Inside a single assignment the right-hand side is analysed before
the store, so the in-place stencil \texttt{f[0,0]}~$\leftarrow$~\texttt{f[1,0]} is accepted, although its
result depends on the loop order. Keep sources and destinations in separate fields.

\begin{keyidea}[Summary]
  \begin{itemize}[leftmargin=1.2em]
    \item \code{create_kernel(assignments, config)} turns the symbolic description into a
          kernel; the \code{CreateKernelConfig} decides target, types and optimizations.
    \item The iteration space is set by ghost layers (inferred by default), a slice, or a
          sparse index field; the layout of the fields fixes the loop order.
    \item On vector CPU targets, enable \code{cpu.vectorize} together with
          \code{assume_inner_stride_one} and pass contiguous arrays; add
          \code{cpu.openmp} for threads.
    \item GPU kernels map cells to threads according to \code{gpu.indexing_scheme}; the
          launch configuration is computed at call time.
  \end{itemize}
\end{keyidea}
